% TOP_PROBLEM: {"id": 4700027, "title": "A moments problem II", "category_id": 4, "category": {"id": 4, "name": "algebra", "display_name": "Algebra", "description": "Group theory, ring theory, field theory, and algebraic structures.", "slug": "algebra", "order_index": 4, "created_at": "2026-07-31T15:26:25.671Z"}, "status": "open", "problem_number": "AMR-046-0027", "set_id": 13}
% FIRST_POSTED: 2026-10-02
% ENTRY_KIND: statement_only
% UnsolvedMath ID 4700027; source catalogue code AMR-046-0027
% !TeX program = lualatex
% Definition and sources only; this file is not an LLM solution attempt.
\documentclass[11pt,a4paper]{article}
\usepackage[margin=25mm]{geometry}
\usepackage{fontspec}
\setmainfont{lmroman10-regular.otf}[BoldFont=lmroman10-bold.otf,
 ItalicFont=lmroman10-italic.otf,BoldItalicFont=lmroman10-bolditalic.otf]
\usepackage{amsmath,amssymb,mathtools}
\usepackage{unicode-math}
\setmathfont{latinmodern-math.otf}
\AtBeginDocument{\let\setminus\smallsetminus}
\usepackage{xurl,hyperref}
\hypersetup{colorlinks=true,urlcolor=blue}
\setcounter{secnumdepth}{0}
\setlength{\parindent}{0pt}
\setlength{\parskip}{5pt}
\setlength{\emergencystretch}{3em}
\begin{document}
\section*{A moments problem II}\label{4700027}
{\small UnsolvedMath ID 4700027 · AMR-046-0027 · Algebra · AMR Open Problem Lists}
\subsection{Definitions and mathematical statement}
For a positive integer $k$, let $\mathcal P_k$ consist of complex
polynomials in one variable with at most $k$ nonzero monomial terms,
including the zero polynomial. A nonzero member can be written
\[
 f(x)=\sum_{j=1}^r c_jx^{d_j},\qquad
 r\le k,\quad 0\le d_1<\cdots<d_r,\quad c_j\in\mathbb C\setminus\{0\}.
\]
The exponents are integers and have no prescribed upper bound. Define
$M_m(f)=\int_0^1 f(x)^m\,dx$ for $m\ge1$, without conjugating $f$.
The question is whether there exists a finite integer $N(k)$ such that
\[
 f\in\mathcal P_k,\quad
 M_1(f)=\cdots=M_{N(k)}(f)=0
 \quad\Longrightarrow\quad f\equiv0.
\]
If such integers exist, determine the least possible values or give useful
explicit bounds in terms of $k$ alone.

The essential uniformity is over all exponent choices as well as all
complex coefficients. A bound depending on the largest degree would not
answer the question. Using ``at most $k$'' includes lower-term degenerations
and gives the monotone version of the source's $k$-monomial question.
The known fact that vanishing of every moment forces a one-variable
polynomial to be zero does not by itself supply a finite bound independent
of its degree. For real coefficients the second moment suffices; the
problem specifically includes complex cancellation.
\subsection{Short English statement}
Can finitely many moments detect every nonzero complex polynomial with $k$ terms, independently of its degree?
\subsection{Sources}
\begin{itemize}
\item[S1] ULAM.AI and the UnsolvedMath Contributors, supplied problems.json, record 4700027 (AMR-046-0027), AMR Open Problem Lists. Catalogue identity and source transcription; CC BY 4.0.
\url{https://huggingface.co/datasets/ulamai/UnsolvedMath}.
\item[S2] Gasull - Some open problems in low dimensional dynamical systems (2020), Problem environment 27: A moments problem II. Original source indexed by AMR-046-0027.
\url{https://arxiv.org/abs/2012.02524}.
\end{itemize}
\end{document}
